% Item compendium template for D&D 5e items with type/rarity grouping
% Uses dndbook document class with structured item organization

% Base LaTeX template for D&D 5e documents using DND-5e-LaTeX-Template
% This template provides the foundation for all D&D-style documents with structured content

\documentclass[letterpaper, 11pt, bg=full, twocolumn, stats=modern]{dndbook}

% DND-5e-LaTeX-Template provides most styling automatically
% Only include additional packages if needed

% Additional packages for enhanced functionality
\usepackage{graphicx}
\usepackage{wrapfig}
\usepackage{hyperref}
\usepackage{bookmark}
\usepackage{makeidx}

% Additional packages for item compendium layout
\usepackage{multicol}
\usepackage{enumitem}
\usepackage{tcolorbox}
\usepackage{titlesec}
\usepackage{tabularx}
\usepackage{array}

% Enhanced list environments for item details
\tcbuselibrary{breakable}

\title{Item Compendium}
\date{\today}

% Enable automatic numbering for chapters only, not sections
\setcounter{secnumdepth}{0}

% Disable text justification but preserve paragraph indentation
\raggedright
\raggedbottom
% Restore paragraph indentation after \raggedright resets it
\setlength{\parindent}{0.8em}

% Hyperlink configuration
\hypersetup{
    colorlinks=true,
    linkcolor=black,
    urlcolor=black,
    citecolor=black,
    filecolor=black,
    bookmarks=true,
    bookmarksopen=true,
    bookmarksnumbered=true,
    pdftitle={Item Compendium},
    pdfsubject={D&D 5e Content},
    pdfkeywords={D&D, 5e, RPG, tabletop, }
}

% Enhanced section styling for item compendium
% DND-5e-LaTeX-Template provides all styling automatically

% Configure for item compendium structure
\setcounter{secnumdepth}{0}  % No section numbers for cleaner look
\setcounter{tocdepth}{2}     % Include subsections in table of contents

% Studiorum Image Helpers — simple, explicit commands
% Load required packages locally so templates don't need to manage them.
\RequirePackage{graphicx}
\RequirePackage{xparse}
\RequirePackage{caption}

% Inline (non-floating) image block
% Usage: \StudiorumImageInline[<width fraction>]{file}{caption}[label]
\NewDocumentCommand{\StudiorumImageInline}{ s o m m o }{%%
  \par\medskip
  \noindent\centering
  \IfNoValueTF{#2}{%%
    \includegraphics[width=\linewidth]{#3}%%
  }{%%
    \includegraphics[width=#2\linewidth]{#3}%%
  }%%
  \IfBooleanF{#1}{%%
    \par\smallskip
    \captionof{figure}{#4}%%
    \IfNoValueF{#5}{\label{#5}}%%
  }%%
  \par\medskip
}

% Standard floating figure
% Usage: \StudiorumImageFloat[<width fraction>]{file}{caption}[label]
\NewDocumentCommand{\StudiorumImageFloat}{ s o m m o }{%%
  
\begin{figure}[htbp]
    \centering
    \IfNoValueTF{#2}{%%
      \includegraphics[width=\linewidth]{#3}%%
    }{%%
      \includegraphics[width=#2\linewidth]{#3}%%
    }%%
    \IfBooleanF{#1}{%%
      \caption{#4}%%
      \IfNoValueF{#5}{\label{#5}}%%
    }%%
  \end{figure}
%%
}

% Wide (two-column spanning when applicable)
% Usage: \StudiorumImageWide[<width fraction>]{file}{caption}[label]
\NewDocumentCommand{\StudiorumImageWide}{ s o m m o }{%%
  \ifdim\columnwidth=\textwidth
    
\begin{figure}[htbp]
      \centering
      \IfNoValueTF{#2}{%%
        \includegraphics[width=\linewidth]{#3}%%
      }{%%
        \includegraphics[width=#2\linewidth]{#3}%%
      }%%
      \IfBooleanF{#1}{%%
        \caption{#4}%%
        \IfNoValueF{#5}{\label{#5}}%%
      }%%
    \end{figure}
%%
  \else
    
\begin{figure*}[tp]
      \centering
      \IfNoValueTF{#2}{%%
        \includegraphics[width=\textwidth]{#3}%%
      }{%%
        \includegraphics[width=#2\textwidth]{#3}%%
      }%%
      \IfBooleanF{#1}{%%
        \caption{#4}%%
        \IfNoValueF{#5}{\label{#5}}%%
      }%%
    \end{figure*}
%%
  \fi
}

% Full-page (dedicated page; fits within text block)
% Usage: \StudiorumImageFullpage[<width fraction>]{file}{caption}[label]
\NewDocumentCommand{\StudiorumImageFullpage}{ s o m m o }{%%
  \clearpage
  
\begin{figure}[p]
    \centering
    \IfNoValueTF{#2}{%%
      \includegraphics[width=\textwidth,height=\textheight,keepaspectratio]{#3}%%
    }{%%
      \includegraphics[width=#2\textwidth,height=\textheight,keepaspectratio]{#3}%%
    }%%
    \IfBooleanF{#1}{%%
      \caption{#4}%%
      \IfNoValueF{#5}{\label{#5}}%%
    }%%
  \end{figure}

  \clearpage
}

% Full-page bleed (fills the physical page; may distort aspect ratio)
% Usage: \StudiorumImageFullpageBleed{file}{caption}[label]
\NewDocumentCommand{\StudiorumImageFullpageBleed}{ s o m m o }{%%
  \clearpage
  
\begin{figure}[p]
    \centering
    \includegraphics[width=\paperwidth,height=\paperheight]{#3}%%
    \IfBooleanF{#1}{%%
      \caption{#4}%%
      \IfNoValueF{#5}{\label{#5}}%%
    }%%
  \end{figure}

  \clearpage
}

% No-caption convenience aliases
\NewDocumentCommand{\StudiorumImageInlineNoCaption}{ o m }{\StudiorumImageInline*[#1]{#2}{}}
\NewDocumentCommand{\StudiorumImageFloatNoCaption}{ o m }{\StudiorumImageFloat*[#1]{#2}{}}
\NewDocumentCommand{\StudiorumImageWideNoCaption}{ o m }{\StudiorumImageWide*[#1]{#2}{}}
\NewDocumentCommand{\StudiorumImageFullpageNoCaption}{ o m }{\StudiorumImageFullpage*[#1]{#2}{}}
\NewDocumentCommand{\StudiorumImageFullpageBleedNoCaption}{ m }{\StudiorumImageFullpageBleed*{#1}{}}

\begin{document}

\tableofcontents
\clearpage

\section{None}

\vspace{0.5em}

\addcontentsline{toc}{subsection}{Amulet of Health}

% Item header with name and metadata
\DndItemHeader{Amulet of Health}{Wondrous item, None, rare (requires attunement)}

% Item statistics for weapons, armor, etc.

% Item properties integrated into description

% Item description
Your Constitution score is 19 while you wear this amulet. It has no effect on you if your Constitution score is already 19 or higher without it.

\vspace{0.8em}
\addcontentsline{toc}{subsection}{Bag of Holding}

% Item header with name and metadata
\DndItemHeader{Bag of Holding}{Wondrous item, None, uncommon}

% Item statistics for weapons, armor, etc.

% Item properties integrated into description

% Item description
This bag has an interior space considerably larger than its outside dimensions, roughly 2 feet in diameter at the mouth and 4 feet deep. The bag can hold up to 500 pounds, not exceeding a volume of 64 cubic feet. The bag weighs 15 pounds, regardless of its contents. Retrieving an item from the bag requires an action.

If the bag is overloaded, pierced, or torn, it ruptures and is destroyed, and its contents are scattered in the Astral Plane. If the bag is turned inside out, its contents spill forth, unharmed, but the bag must be put right before it can be used again. Breathing creatures inside the bag can survive up to a number of minutes equal to 10 divided by the number of creatures (minimum 1 minute), after which time they begin to suffocate.

Placing a bag of holding inside an extradimensional space created by a \textit{Heward's handy haversack}, \textit{portable hole}, or similar item instantly destroys both items and opens a gate to the Astral Plane. The gate originates where the one item was placed inside the other. Any creature within 10 feet of the gate is sucked through it to a random location on the Astral Plane. The gate then closes. The gate is one-way only and can't be reopened.

\vspace{1.5em}

\section{Melee Weapon}

\vspace{0.5em}

\addcontentsline{toc}{subsection}{Longsword}

% Item header with name and metadata
\DndItemHeader{Longsword}{Melee weapon}

% Item statistics for weapons, armor, etc.
\textit{15 gp}

% Item properties integrated into description

% Item description

\end{document}